\chapter{ความรู้เบื้องต้นเกี่ยวกับระบบดิจิทัล}

ระบบวงจรทางอิเล็คทรอนิคในปัจจุบันสามารถแบ่งออกเป็น 2 ระบบได้แก่ ระบบ ดิจิตอล (digital) และ ระบบอนาล๊อค (analog) ซึ่งระบบดิจิตอลจะแสดงถึงปริมาณสัญญาณชนิดไม่ต่อเนื่อง (discrete signal) ส่วนระบบแบบอนาล๊อคจะแสดงถึงปริมาณสัญญาณแบบต่อเนื่อง (continuous signal) โดยที่ปริมาณสัญญาณแบบต่อเนื่องนี้ จะเกิดขึ้นทั่วไปจากธรรมชาติ เช่น เสียง แสง แรงดันไฟฟ้า เป็นต้น เราสามารถนำปริมาณสัญญาณแบบต่อเนื่องนี้ มาเปลี่ยนรูปให้อยู่ในรูปแบบสัญญาณแบบไม่ต่อเนื่องหรือระบบดิจิตอลได้ ซึ่งในระบบการทำงานของคอมพิวเตอร์ยุคปัจจุบันจะทำงานแบบระบบดิจิตอลเป็นหลัก การแปลงปริมาณสัญญาณแบบ อนาล๊อคเป็นดิจิตอลนั้น จำเป็นต้องเข้าใจข้อแตกต่างระหว่างระบบทั้งสองแบบ ตัวอย่างเช่น ปริมาณการเปลี่ยนแปลงอุณหภูมิในอากาศภายใน 1 วัน เราสามารถแสดงเป็นกราฟความสัมพันธ์ ระหว่างอุณหภูมิตามช่วงเวลาต่างๆ ได้ดังภาพที่ \ref{fig:continueSignal}

\begin{figure}[!htb]
    \centering
    \includegraphics[height=7cm]{image001.png}
    \caption{กราฟแสดงความสัมพันธ์ของการเปลี่ยนแปลงอุณหภูมิตามช่วงเวลาต่างๆ ในลักษณะสัญญาณแบบต่อเนื่อง}
    \label{fig:continueSignal}
\end{figure}

จากภาพจะเห็นว่าอุณหภูมิมีการเปลี่ยนแปลงตลอดเวลา เมื่อลองพิจารณาช่วงการเปลี่ยนแปลงของอุณภูมิระหว่าง 58 องศา ถึง 59 องศา จะเห็นว่าจำนวนตัวเลขที่อยู่ในช่วงระหว่าง 58-59 องศา มีมากถึง infinity จำนวน ตัวอย่างเช่น 58.00001, 58.000002, …. เป็นต้น ซึ่งหากจำนวนตัวเลขเหล่านี้มีจำนวนมากส่งผลให้ กราฟการเปลี่ยนแปลงมีความต่อเนื่องและละเอียดมากยิ่งขึ้น
	สำหรับในระบบดิจิตอลแล้ว เราไม่สามารถทำการเก็บค่าจำนวนที่มีความละเอียดสูง ที่อาจจะเข้าใกล้ infinity ได้ดังนั้น เราจำเป็นต้องมการแปลงปริมาณสัญญาณที่เป็น อนาล๊อค ให้กลายเป็นปริมาณสัญญาณทางดิจิตอล โดยอุปกรณ์ที่ทำหน้าที่แปลงสัญญาณนี้จะเรียกว่า ตัวแปลงสัญญาณอนาล็อคเป็นดิจิตอล (Analog to digital converter :ADC) หรือในทางกลับกับ อุปกรณ์ที่ทำหน้าที่แปลงสัญญาณจากดิจิตอลเป็นอนาล๊อค จะเรียกว่า ตัวแปลงดิจิตอลเป็นอนาล๊อค (Digital to analog converter: DAC)

\section{การแปลงสัญญาณอนาล๊อคเป็นดิจิตอล}
การแปลงสัญญาณอนาล๊อคเป็นดิจิตอลนั้นดังที่กล่าวมาข้างต้นจำเป็นอย่างยิ่งที่จะต้องมีอุปกรณ์ในการเปลี่ยนรูปแบบของปริมาณสัญญาณที่ต้องการ ในอดีตอุปกรณ์ทางไฟฟ้า แทบทุกชนิดใช้ระบบแบบอนาล๊อคซึ่งมีปัญญาในด้านการจัดเก็บและการนำมาประมวลผล ตัวอย่างเช่น ระบบเครื่องขยายที่เป็นระบบแบบอนาล๊อค ดังภาพที่  \ref{fig:ampifier} จะรับคลื่นเสียงผ่านไมโครโฟน ซึ่งทำหน้าที่แปลงสัญญาณเสียงที่เป็นสัญญาณอนาล๊อคให้กลางเป็นสัญญาณไฟฟ้าแบบอนาล๊อค
	\begin{figure}
	\centering
	\includegraphics[height=5cm]{system01.jpg}
	\caption{ระดับสัญญาณไฟฟ้าและสัญญาณลอจิก}
	\label{fig:ampifier}
\end{figure}
 ซึ่งสัญญาณไฟฟ้าที่ได้จากไมโครโฟนนั้นเป็นสัญญาณขนาดเล็ก (small amplitude) จึงจำเป็นต้องทำการขยายสัญญาณให้มีขนาดที่สูงขึ้น (high amplitude) แต่เมื่อมีการขยายสัญญาณ ส่งผลให้เกิดสัญญาณรบกวนมากขึ้นเช่นกันอีกทั้งยังพบปัญหาในการจัดเก็บและค้นหาข้อมูลอีกด้วย  ดังนั้นระบบสัญญาณแบบดิจิตอลจึงเข้ามาแทนที่ระบบสัญญาณแบบอนาล๊อค แต่ข้อจำกัดของระบบสัญญาณแบบดิจิตอลก็คือความละเอียดของการจัดเก็บหรือบันทึกจะไม่สามารถเก็บเก็บความละเอียดได้ทั้งหมดเหมือนกับสัญญาณอนาล๊อค
ตัวอย่าง ให้ทำการพิจารณา ภาพที่  \ref{fig:continueSignal} จะพบว่าทุกช่วงเวลาเมื่อลากเส้นตัดไปจะมีค่าสัญญาณในแกน y เสมอ แต่หากเป็นสัญญาณดิจิตอลซึ่งเป็นสัญญาณไม่ต่อเนื่อง เราต้องทำการระบุว่าจะทำการบันทึกข้อมูลแต่ละช่วงห่างกันเท่าใด ซึ่งระยะห่างของช่วงเวลาที่บันทึกมีผลทำให้ค่าสัญญาณที่แปลงมานั้นมีความละเอียดมากหรือละเอียดน้อยขึ้นกับช่วงเวลาการบันทึกดังกล่าว 

สรุปสัญญาณที่เกิดขึ้นตามธรรมชาติส่วนใหญ่แล้วเป็นสัญญาณอนาล๊อคซึ่งทำการเก็บข้อมูลได้ยากในระบบคอมพิวเตอร์หรือระบบดิจิตอลดังนั้นจึงจำเป็นต้องมีการแปลงสัญญาณให้อยู่ในรูปแบบอื่นๆ ซึ่งการแปลงสัญญาณไปเป็นรูปแบบอื่นๆ จะทำให้เกิดการสูญเสียลักษณะสัญญาณไปไม่สามารถเก็บความละเอียดได้ทั้งหมดขึ้นอยู่กับวิธีการแปลงการเก็บข้อมูลนั้นๆ
\section{ลักษณะสัญญาณดิจิตอล}
ในระบบดิจิตอลจะมีเพียงสัญญานลอจิกหนึ่ง หรือ Logical high และ อีกสัญญาณหนึ่งคือ ลอจิกศูนย์ หรือ Logical Low ซึ่งสัญญาณ 1 สัญญาณจะมีสถาณะได้เพียง 1 ค่าเท่านั้นนั่นคือ ไม่เป็น ลอจิกศูนย์ หรือไม่ก็ ลอจิกหนึ่ง เราจะเรียกสํญญาณ 1 สัญญาณนี้ว่า 1 บิต ซึ่งเป็นหน่วยที่เล็กที่สุดของระบบดิจิตอล 
ค่าสัญญาณที่เป็นลอจิก 1 หรือ ลอจิก 0 จะใช้ค่าระดับแรงดันไฟฟ้า โดยทั่วไปแล้วอุปกรณ์อิเล็คทรอนิคต่างๆ จะนิยามให้ ลอจิก 1 มีระดับแรงดันไฟฟ้าอยู่ที่ 5 โวลต์ และลอจิก 0 อยู่ที่ระดับแรงดันไฟฟ้าที่ 0 โวลต์ ทั้งนี้ขึ้นอยู่กับการนิยามของอุปกรณ์อิเล็คทรอนิคนั้นๆ ตัวอย่างเช่น
อุปกรณ์บางชนิดอาจนิยามให้ลอจิก 1 ใช้ระดับแรงดันไฟฟ้าอยู่ที่ 12 โวลต์ และ ลอจิก 0 อยู่ที่ -12 โวลต์ หรือ อาจจะใช้ ลอจิก 1 อยู่ที่ 3 โวลต์และลอจิก 0 อยู่ที่ 0 โวลต์ ซึ่งสามารถทำได้เช่นกันขึ้นอยู่กับอุปกรณ์อิเล็คโทรนิคนั้นๆ
	เนื่องจากแรงดันไฟฟ้าเป็นสัญญาณอนาล็อค การที่จะเปรียบเทียบระดับแรงดันไฟฟ้าให้เป็นลอจิก 1 หรือ 0 ผู้พัฒนาอุปกรณ์ จะกำหนดช่วงชองค่าแรงดันไฟฟ้าเพื่อสำหรับเปรียบเทียบให้เป็นสัญญาณลอจิก 1 หรือ 0 ดังภาพที่  \ref{fig:logicLevel}
	
	\begin{figure}
	    \centering
	    \includegraphics[height=5cm]{image003.png}
	    \captionsource{ระดับสัญญาณไฟฟ้าและสัญญาณลอจิก}{\cite{bri_logic_2020}}
	    \label{fig:logicLevel}
	\end{figure}
	
		จากภาพที่  \ref{fig:logicLevel} จะแสดง $V_{IL}$ แสดงถึงระดับสัญญาณแรงดันไฟฟ้าของสัญญาณนำเข้า เมื่อลอจิกเป็น 0 และลอจิกจะเป็น 1 เมื่อระดับแรงดันไฟฟ้าของสัญญาณนำเข้ามากกว่า $V_{IH}$ ทางด้านของสัญญาณส่งออกก็เช่นกัน เมื่ออุปกรณ์ส่งค่าลอจิก 1 นั่นหมายถึงอุปกรณ์อิเล็คทรอนิคจะส่งระดับแรงดันไฟฟ้าที่ $V_{OH}$ ขึ้นไป และหากส่งลอจิก 0 จะส่งระดับแรงดันไฟฟ้าที่ $V_{OL}$ ลงมาจนถึง 0 โวลต์
       
       สรุป ระบบสัญญาณดิจิตอลเป็นสัญญาณที่ไม่ต่อเนื่อง โดยจะมี สองสถาณะนั่นคือลอจิกหนึ่ง หรือ ลอจิกศูนย์ โดยที่ใน 1 ช่วงเวลาสามารถเป็นได้เพียง 1 สถาณะเท่านั้น และเราจะเสียกสัญญาณลอจิก 1 สัญญาณนี้ว่า 1 บิต ซึ่งเป็นหน่วยที่เล็กที่สุดของระบบดิจิตอล
       
\section{ระบบสัญญาณนาฬิกาและสัญญาณพัลส์}
สำหรับการทำงานร่วมกันของอุปกรณ์ทางอิเล็คทรอนิกส์ หลายๆ ชนิดนั้น อุปกรณ์เหล่านี้จะต้องมีตัวบอกจังหวะของการทำงาน เราจะเรียกว่า สัญญาณนาฬิกา (Clock) ซึ่งทำหน้าที่ในการกระตุ้นการทำงาน หรือ คอยควบคุมจังหวะเวลาการทำงานของอุปกรณ์อิเล็คทรอนิคต่างๆ โดยสัญญาณนาฬิกาจะมีค่าลอจิกสลับกันระหว่าง ลอจิก 0 และ ลอจิก 1 วนไปเรื่อยๆ ทั้งนี้คาบของการเปลี่ยนแปลงของลอจิกจะคงที่เสมอ ซึ่งทำให้ความถี่ของสัญญาณนาฬิกาคงที่ค่าหนึ่ง  ตัวอย่างของสัญญาณนาฬิกาแสดงดังภาพที่ \ref{fig:clocksignal}

\begin{figure}[!htb]
    \centering
    \includegraphics[height=5cm]{ClockCycle.jpg}
    \caption{ตัวอย่างของสัญญาณนาฬิกา}
    \label{fig:clocksignal}
\end{figure}

จากภาพที่  \ref{fig:clocksignal} จะเห็นว่า การช่วงสัญญาณนาฬิกาที่เป็นลอจิก 1 และ ลอจิก 0 เราจะเรียกว่า หนึ่งคาบสัญญาณ (Time Period) ซึ่งสามารถหาความถี่ของสัญญาณนาฬิกาได้จากสมการ

\begin{equation}
    f = \frac{1}{T}, \text{ and } T=\frac{1}{f},
    \label{eq:feqcal}
\end{equation}

โดยที่ $f$ แทนค่าความถี่ (Frequency) และ $T$  แทนคาบของสัญญาณนาฬิกา ลักษณะของสัญญาณนาฬิกานั้นจะให้ช่วงของลอจิก 1 และ 0 ที่เท่ากันซึ่งทำให้ลักษณะสัญญาณคงที่เสมอ (Periodic) แต่มีสัญญาณอีกชนิดหนึ่งนั่นคือสัญญาณพัลส์จะเป็นสัญญาณที่ส่งออกมาช่วงหนึ่งๆ เพื่อกระตุ้นการทำงาน อาจจะมีลักษณะสัญญาณที่คงที่เสมอหรือไม่ก็ได้ (Non-periodic) สัญญาณพัลส์มีด้วยกัน 2 ชนิดนั่นคือสัญญาณพัลส์บวก (Positive logic pulse) และ พัลส์ลบ (Negative logic pulse) แสดงดังภาพที่  \ref{fig:plusesignal}

\begin{figure}[!htb]
    \centering
    \includegraphics{image016.jpg}
    \caption{ภาพสัญญาณพัลส์บวก และ พัลส์ลบ}
    \label{fig:plusesignal}
\end{figure}

ลักษณะสัญญาณของสัญญาณพัลส์ในระบบดิจิตอลไม่จำเป็นต้องมีช่วงของลอจิก 1 และ ลอจิก 0 ที่เท่ากันเสมอไป โดยเราสามารถคำนวนอัตราส่วนระหว่างค่าช่วงของสัญญาณลอจิก 1 กับคาบเวลาโดยมีค่าเป็นเปอร์เซ็นต์ เราจะเรียกอัตราส่วนนี้ว่า ดิวตี้ไซเคิล (Duty cycle) ซึ่งค่าดิวตี้ไซเคิล สามารถคำนวณได้ดังสมการต่อไปนี้

\begin{equation}
    Duty\ cycle =\frac{t_w}{T}\times 100 \%
    \label{eq:dutycal}
\end{equation}

ซึ่งอัตราส่วน Duty cycle สามารถแสดงได้ดังรูป \ref{fig:sample dutycycle}
\begin{figure}[!htb]
    \centering
    \includegraphics[width=0.8\textwidth]{image019.png}
    \caption{ตัวอย่างสัญญาณพัลส์ในแต่ละค่า Duty cycle ที่แตกต่างกัน}
    \label{fig:sample dutycycle}
\end{figure}

\begin{exmp}
\normalfont {กำหนดให้สัญญาณพัลส์มีคาบเวลาที่แน่นอนและมีลักษณะดังรูป  จงหาคาบ ความถี่ และ ค่าดิวตี้ไซเคิล กำหนดให้ 1 ช่องมีช่วงเวลาเท่ากับ 20 ns}
\end{exmp}
\begin{figure}[h]
    \centering
    \includegraphics[width=0.8\textwidth]{image021.png}
    \caption{ตัวอย่างภาพสัญญาณพัลส์}
    \label{fig:exam1}
\end{figure}
\par จากภาพที่  \ref{fig:exam1} จะเห็นว่า 1 คาบของสัญญาณกินช่องของเวลาไป 6 ช่องเวลา ดังนั้น 1 คาบใช้เวลาทั้งหมด เท่ากับ $6 \times 20 \ ns = 120 \ ns$   ซึ่งสามารถหาความถี่ได้โดยใช้สมการ \eqref{eq:feqcal} ดังนั้นความถี่มีค่าเท่ากับ $\frac{1}{120\ ns} = \ 8.34\ MHz$ สำหรับค่าดิวตี้ไซเคิล ให้พิจารณาช่วงเวลาที่ สัญญาณมีค่าเป็นลอจิก 1 ซึ่งจะพบมา กินช่วงเวลา 1 ช่วง ซึ่งมีค่าเท่ากับ $20 \ ns$ โดยสามารถคำนวนค่า ดิวตี้ไซเคิลได้จากสมการ \eqref{eq:dutycal} ได้เท่ากับ $dutycycle= \frac{t}{T})\times100 \%=\frac{20}{120}\times100\%=16.6\%$

สรุป ระบบสัญญาณนาฬิกาเป็นสัญญาณที่มีรูปคลื่นที่แน่นอน มีคาบและความถี่ที่ชัดเจน โดยสามารถคำนวนความสัมพันธ์ระหว่างคาบและความถี่ได้จากสมการ \eqref{eq:dutycal} และ \eqref{eq:feqcal} สัญญาณนาฬิกาจะมีช่วงของเวลาที่มีลอจิก 1 และ ลอจิก 0 มีช่วงเวลาที่เท่า ๆ กัน แต่สัญญาณพัลส์จะสามารถปรับแต่งให้มีช่วงเวลาที่เป็นลอจิก 1 และ 0 ที่ไม่เท่ากันได้ ซึ่งอัตราส่วนของช่วงเวลาที่มีสถาณะเป็นลอจิก 1 กับ 0 เราเรียกว่า ค่าดิวตี้ไซเคิ้ล


\section{เครื่องมือวัดสำหรับระบบงานดิจิตอล}

การวัดค่าสัญญาณเป็นสิ่งที่จำเป็นสำหรับในระบบดิจิตอลเนื่องจาก ทำให้เราสามารถทราบถึงลักษณะสัญญาณ เช่น ค่าของสัญญาณ ความถี่ของสัญญาณ และรูปแบบของสัญญาณ สำหรับเครื่องมือวัดในระบบงานดิจิตอล มีเครื่องมือวัดหลากหลายชนิด อันได้แก่

\subsection{มัลติมิเตอร์ดิจิตอล(Digital Multimeters)} 

มัลติมิเตอร์เป็นเครื่องวัดระดับสัญญาณไฟฟ้าที่สามารถวัดปริมาณไฟฟ้าได้ หรือเป็นเครื่องค่าทางไฟฟ้าประเภทอื่นๆ ได้เช่น วัดกระแสไฟฟ้าทั้งกระแสสลับและกระแสตรง แรงดันไฟฟ้าทั้งกระแสสลับและกระแสตรง ความต้านทาน หรือ ทดสอบการทำงานของไดโอดได้ ซึ่งมัลติมิเตอร์มีทั้งแบบที่เป็นการแสดงค่าที่วัดได้ด้วยระบบดิจิตอล และแบบเข็มชี้วัด โดยทั้ง 2 แบบมีข้อดีที่แตกต่างกัน ในมัลติมิเตอร์แบบดิจิตอล สามารถให้ค่าที่เที่ยงตรงและแม่นยำเนื่องจากแสดงค่าออกมาเป็นตัวเลขชัดเจน แต่จะไม่เห็นถึงระดับการเปลี่ยนแปลงของสัญญาณ สำหรับในมัลติมิเตอร์แบบเข็ม จะต้องทำการอ่านค่าจากเข็มที่ชี้ไปที่สเกลวัด ทำให้การอ่านค่ามีความไม่เที่ยงตรงตามผู้ที่ทำการอ่านค่า แต่จะเห็นลักษณะการเปลี่ยนแปลงของสัญญาณไฟฟ้า ตามการแกว่งเข็มของมัลติมิเตอร์ 
ซึ่งรูปมัลติมิเตอร์ทั้งสองแบบแสดงดังภาพที่ \ref{fig:multimeter}

%\begin{figure}
%	\centering
%	\includegraphics[width=0.5\textwidth]{multimeter.jpg}
%	    \caption{ตัวอย่างเครื่องมือวัด มัลติมิเตอร์}
%	    \Source{www.dummy.com}
%	\label{fig:multimeter}
%\end{figure}


 \begin{figure}[H]
	\centering
	\captionsetup[subfigure]{format=plain,justification=centering}
	\begin{subfigure}[b]{0.4\textwidth}
		\includegraphics[height=5cm]{multimeter.jpg}
	\end{subfigure}
	\begin{subfigure}[b]{0.4\textwidth}
		\includegraphics[height=5cm]{Analogmultimeter.jpg}
	\end{subfigure}
	\captionsource{ตัวอย่างเครื่องมือวัดมัลติมิเตอร์แบบดิจิตอลและแบบเข็ม}{\citep{mihai_digital_2021};
		\cite{csaba_analog_2021}}
	\label{fig:multimeter}
\end{figure}


\subsection{ออสซิลโลสโคป(Oscilloscope)}

ออสซิลโลสโคปเป็นเครื่องมือวัดทางไฟฟ้าที่ใช้สำหรับวัดรูปคลื่นสัญญาณไฟฟ้าออกมาเป็นรูปสัญญาณทางจอมอนิเตอร์ ซึ่งสามารถวัดได้ทั้ง ขนาด(Amplitude) และเวลา(Time) โดยที่สามารถบอกให้อยู่ในรูปของค่าความถี่และคาบเวลาได้ โดยออสซิลโลสโคป สามาถวัดค่าสัญญาณทางไฟฟ้าได้ตั้งแต่ 1 ช่องสัญญาณขึ้นไปโดยปกติแล้ว ออสซิลโลสโคปจะมีให้ 2 ช่องสัญญาณ ตัวออสซิลโลสโคปและลักษณะของออสซิลโลสโคปแสดงดังภาพ \ref{fig:osillo}

\begin{figure}[!htb]
	\centering
	\includegraphics[height=5cm]{osiloscope.jpg}
	\captionsource{รูปแสดงเครื่องมือวัดออสซิลโลสโคป}{\cite{lee_isolated_2020}}
	\label{fig:osillo}
\end{figure}
สำหรับการวัดสัญญาณ จะใช้สายเคเบิลที่เรียกว่า  สายโพรบ (Probe) เป็นสายเคเบิลแกนร่วม (Coaxial cable) โดยสายนี้จะเชื่อมต่อกับตัวออสซิลโลสโคป 1 ช่องสัญญาณ และปลายสายอีกด้านจะมี 2 ขั้ว ขั้วหนึ่งทำหน้าที่จับสัญญาณอ้างอิง ในที่นี้เราจะใช้สัญญาณการ์ว (Ground) เป็นสัญญาณอ้างอิง และอีกขั้วหนึ่งจะทำหน้าที่จับสัญญาณเทียบกับสัญญาณอ้างอิง ซึ่งสายโพรบนี้แสดงภาพตัวอย่างดังภาพที่ \ref{fig:probe}

\begin{figure}[!htb]
	\centering
	\includegraphics[height=5cm]{probe.jpg}
 	\captionsource{รูปตัวอย่างสายสัญญาณโพรป}{\cite{david_basics_2020}}
	\label{fig:probe}
\end{figure}



\subsection{ลอจิกอนาไลเซอร์ (logic analyzer) }
ลอจิกอนาไลเซอร์ เปนเครื่องมือที่ใชตรวจดูสัญญาณที่เกิดขึ้นในระบบ ดิจิตอล ส่วนใหญ่มักใช้ตรวจจับการทำงานของการส่งข้อมูลพร้อมๆ กันในไมโครโพรเซสเซอร์ (microprocessor) โดยที่ ลอจิกอนาไลเซอร์นั้นเน้นการดูรูปสัญญาณที่เป็นสัญญาณทางดิจิตอลอย่างเดียว ซึ่งแตกต่างกับ ออสซิโลสโคปที่สามารถดูรูปสัญญาณได้ทั้งสัญญาณอนาล็อค และ ดิจิตอล เนื่องจาก ลอจิกอนาไลเซอร์เป็นเครื่องมือที่ใช้ตรวจจับสัญญาณดิจิตอลเพียงอย่างเดียวดังนั้นจึงทำการเพิ่มความสามารถเพื่อให้ลอจิกอนาไลเซอร์สามารถจับสัญญาณได้หลาย ๆ ช่องสัญญาณพร้อม ๆ กัน ดังนั้นลอจิกอนาไลเซอร์มุ้งเน้นไปที่การจับสัญญาณที่ให้ความแม่นยำในเรื่องของเวลามากกว่า การวัดระดับของแรงดันสัญญาณ

ลอจิกอนาไลเซอร์ สามารถบันทึกขอมูลเพื่อการวิเคราะหข้อมูลทางเวลา (timing
analysis) หรือเพื่อวิเคราะหทางสถานะ (state analysis)
การวิเคราะห์ข้อมูลทางเวลา จะใช้สัญญาณนาฬิกาของ ลอจิกอนาไลเซอร์ เป็นสัญญาณอ้างอิง แต่หากใช้สัญญาณนาฬิกาจากระบบภายนอกที่กำลังตรวจสอบจะเป็นการบันทึกข้อมูลเพื่อใช้สำหรับเป็นการวิเคราะห์ทางสถานะ

เครื่องมือลอจิกอนาไลเซอร์ สามารถบันทึกข้อมูลผ่านทางหน่วยความจำแบบพกพาได้หลายชนิด เช่น Flash drive หรือ SD card ขึ้นอยู่กับผู้ผลิตที่จะออกแบบความสามารถในการเก็บข้อมูลของลอจิอนาไลเซอร์ ลักษณะของเครื่องมือ ลอจิกอนาไลเซอร์มีลักษณะดังภาพที่  \ref{fig:logicanalizer}

\subsection{เครื่องกำเนิดสัญญาณ (The Function Generator)}

เครื่องกำเนิดสัญญาณ เป็นอุปกรณ์ที่จำเป็นอย่างมากในการสร้างสัญญาณ ในรูปแบบต่างๆ เช่น แบบคลื่นพัลส์ (Pulse wave), แบบคลื่นไซส์ (Sin wave) หรือแบบคลื่นสามเหลี่ยม (triangular wave) เป็นต้น ซึ่งใช้เป็นตัวกำหนดสัญญาณและใช้สำหรับตรวจเช็คการทำงานของระบบดิจิตอล  วัดเปรียบเทียบค่า หรือใช้อ้างอิง นำไปใช้งานในวงจรหรือเครื่องมืออุปกรณ์ต่างๆ ทางด้านไฟฟ้าและอิเล็กทรอนิกส์ นับว่ามีความสัญคัญอย่างมากในการ ตรวจวัดและทดสอบสัญญาณ
 โดยลักษณะของเครื่องกำเนิดสัญญาณมีลักษณะดังรูป


\section{สรุป}
สรุปสัญญาณที่เกิดขึ้นตามธรรมชาติส่วนใหญ่แล้วเป็นสัญญาณอนาล๊อคซึ่งทำการเก็บข้อมูลได้ยากในระบบคอมพิวเตอร์หรือระบบดิจิตอลดังนั้นจึงจำเป็นต้องมีการแปลงสัญญาณให้อยู่ในรูปแบบอื่นๆ ซึ่งการแปลงสัญญาณไปเป็นรูปแบบอื่นๆ จะทำให้เกิดการสูญเสียลักษณะสัญญาณไปไม่สามารถเก็บความละเอียดได้ทั้งหมดขึ้นอยู่กับวิธีการแปลงการเก็บข้อมูลนั้นๆ

ระบบสัญญาณดิจิตอลเป็นสัญญาณที่ไม่ต่อเนื่อง โดยจะมี สองสถาณะนั่นคือลอจิกหนึ่ง หรือ ลอจิกศูนย์ โดยที่ใน 1 ช่วงเวลาสามารถเป็นได้เพียง 1 สถาณะเท่านั้น และเราจะเสียกสัญญาณลอจิก 1 สัญญาณนี้ว่า 1 บิต ซึ่งเป็นหน่วยที่เล็กที่สุดของระบบดิจิตอล


ระบบสัญญาณนาฬิกาเป็นสัญญาณที่มีรูปคลื่นที่แน่นอน มีคาบและความถี่ที่ชัดเจน โดยสามารถคำนวนความสัมพันธ์ระหว่างคาบและความถี่ได้จากสมการ \eqref{eq:dutycal} และ \eqref{eq:feqcal} สัญญาณนาฬิกาจะมีช่วงของเวลาที่มีลอจิก 1 และ ลอจิก 0 มีช่วงเวลาที่เท่า ๆ กัน แต่สัญญาณพัลส์จะสามารถปรับแต่งให้มีช่วงเวลาที่เป็นลอจิก 1 และ 0 ที่ไม่เท่ากันได้ ซึ่งอัตราส่วนของช่วงเวลาที่มีสถาณะเป็นลอจิก 1 กับ 0 เราเรียกว่า ค่าดิวตี้ไซเคิ้ล
\begin{figure}
	\centering
	\fbox{\includegraphics[height=7cm]{logicanalyzer.jpg}}
	\captionsource{ลักษณะของลอจิกอนาไลเซอร์}{\cite{herres_logic_2018}}
	\label{fig:logicanalizer}
\end{figure}

\newpage
\QAChapter
\begin{enumerate}
	\item จงอธิบายความแตกต่างของสัญญาณอนาลอคและสัญญาณดิจิตอล 
	\item จงยกตัวอย่างสัญญาณอนาล๊อค และ สัญญาณดิจิตอลมาอย่างละ 5 ตัวอย่างสัญญาณ
	\item ในทางอุปกรณ์อิเล็คโทรนิคนิยมใช้ ตัวชี้วัดใดในการบอกสัญญาณลอจิก 1 และ ลอจิก 0 ในระบบดิจิตอล
	\item จงอธิบายหน้าที่ของสัญญาณนาฬิกาในระบบดิจิตอลทำหน้าที่อะไร
	\item จงหาความถี่ของสัญญาณนี้ เมื่อกำหนดให้ค่าดิวตี้ไซเคิล มีอัตราส่วน 20\% ละ 1 คาบสัญญาณ กินเวลา 50 ms 
	\item จงอธิบายลักษณะการใช้งานของเครื่องมือวัดประเภทมัลติมิเตอร์
	\item จงอธิบายหน้าที่ของเครื่องมือวัดประเภทออซิโลสโคป
	\item จงอธิบายข้อแตกต่างระหว่างเครื่องมือวัดประเภทออซิโลสโคปกับ ลอจิกอนาไลเซอร์
\end{enumerate}
